%%FILE VII-b
%%APPEND title="The universal property includes natural transformations"
Four facts fix what these path categories are known to compute. (1) The set of components of the category of paths from $a$ to $b$, with arbitrary vertical arrows, is the set of morphisms from $a$ to $b$ in the localization of $M$ by $\Sigma$; in the ladder argument no vertical arrow is inverted. (2) Let $J$ be the inclusion of the subcategory of the same paths with vertical arrows in $\Sigma$, the hammocks of Dwyer and Kan. It is not an isomorphism of categories, even for a saturated $\Sigma$: for $M=[1]$ with only identities marked the full category has more arrows. (3) $J$ is a weak equivalence when $\Sigma$ consists of isomorphisms, and it is the identity when every arrow is marked. (4) In general the full path space maps to the mapping space of the simplicial localization, and the composite with $J$ is the comparison of Dwyer and Kan; so $J$ is injective on all homotopy groups and bijective on components. Whether $J$ is always a weak equivalence, that is, whether the path categories of this chapter have the homotopy type of the mapping spaces of the localization, is not known; the first undecided case is the ordered set $0<1<2$ with the arrow $0\to2$ marked. Of these, (4) was followed in outline only.

%%APPEND title="The cofinality condition that controls total categories"
The exact condition is the following. For a functor $u:J\to T$, the induced functor between the total categories of $F$ and of its restriction to $J$ is a weak equivalence for every strict functor $F:T^{\mathrm{op}}\to\mathrm{Cat}$ if and only if every comma category $u/t$ is weakly contractible. Sufficiency is Theorem~A applied over $T$; necessity is seen on the representable functors. Cofinality of the image in the preorder of types, even with homotopic reindexing maps, is not enough, as the example at ms.~109 shows. For a filtered index category the total category is weakly equivalent to the colimit in $\mathrm{Cat}$, with no hypothesis on the transition functors.

%%APPEND title="Independence of the sequential homotopy type"
The comparisons can be made canonical, for every admissible $\Sigma$. For two types $\sigma,\tau$ let $P(\tau,\sigma)$ be the category of the increasing maps of types that preserve end points. It is empty or contractible (commentary at ms.~86), and these categories compose strictly. They act on the path categories, $P(\tau,\sigma)\times F_\sigma\to F_\tau$, and through this action the diagram of path categories descends coherently to the preorder of types; the category of all types and every sequence of types cofinal in the preorder are homotopy final there. It follows that the functor from the total category of a cofinal sequence to the total category over all types is a weak equivalence, that maps of sequences induce weak equivalences, and that these comparisons are transitive. The argument was followed in outline; the passage through the homotopy coherent nerve is not unfolded here. It replaces the argument through ind-objects of $\mathrm{Cat}\,\mathrm{hot}$, which cannot give this (commentary on Chapter~V, ms.~49).

%%FILE VII-c
%%BLOCK after="A comparison that fails under the general hypotheses"
\begin{dcomment}{development}{p121}{The padding functors are weak equivalences and are not invertible modulo natural transformations}
The question of this page has two answers, according to the category in which it is asked. For the path categories themselves (the geometric case, as opposed to a general fibred category over the types), the padding functors that extend a path by constant tails, and the comparison functors between the sequential models and the category over all types, are weak equivalences for every admissible $\Sigma$, and the composition of paths they define is associative and unital in $\mathsf{Hot}$, independently of the block of types used; this follows from the action of the categories of maps of types described in the commentary at ms.~117, and was followed in outline. In $\mathrm{Cat}\,\mathrm{hot}$ the answer is negative.

\textbf{Proposition.}\ours
\textit{Let $X$ be the ordered set with four elements $0,1,2,3$ and the relations $0<1$, $0<3$, $2<1$, $2<3$, with all arrows marked. Let $\mathcal P$ be the category of the functors from the integer fence $J$ (arrows from each even integer to its two neighbours) to $X$ that are constant with value $0$ far to the left and far to the right, with the natural transformations that are identities on both tails, and $\mathcal L\subset\mathcal P$ the full subcategory of the paths constant with value $0$ on the nonnegative integers. The inclusion $q:\mathcal L\to\mathcal P$ is a weak equivalence, and there is no functor $r:\mathcal P\to\mathcal L$ with $qr$ connected to the identity of $\mathcal P$ by a zigzag of natural transformations.}

\textit{Proof.}
That $q$ is a weak equivalence is the general statement above. Let $\widetilde X$ be the integer fence itself; reduction modulo $4$ is a functor $\widetilde X\to X$ along which every arrow of $X$ lifts uniquely once the lift of its source, or of its target, is given. Hence every object $c$ of $\mathcal P$ has a unique lift $\widetilde c:J\to\widetilde X$ with value $0$ far to the left, and a morphism $c\to c'$ of $\mathcal P$ lifts uniquely to a natural transformation $\widetilde c\to\widetilde c'$ that is the identity far to the left. Far to the right this transformation is an identity as well, so the constant value of $\widetilde c$ far to the right is the same for two objects joined by a morphism.

Suppose $r$ exists, with a zigzag of $k$ natural transformations between the identity and $qr$. Let $c$ be an object whose lift has value $0$ far to the right. Evaluating the zigzag at $c$ gives a zigzag of $k$ morphisms of $\mathcal P$ from $c$ to $qr(c)$; all the objects in it have lifts with value $0$ far to the right. The lift of $qr(c)$ is constant on the nonnegative integers, with value $0$. Evaluating the lifted zigzag at an integer $i\geq0$ gives a zigzag of $k$ arrows of $\widetilde X$ from $\widetilde c(i)$ to $0$; an arrow of $\widetilde X$ joins integers at distance at most one, so $|\widetilde c(i)|\leq k$. But for $N\geq1$ the map $\widetilde c_N:J\to\widetilde X$ with $\widetilde c_N(i)=0$ for $i\leq0$ and for $i\geq8N$, $\widetilde c_N(i)=i$ for $0\leq i\leq4N$ and $\widetilde c_N(i)=8N-i$ for $4N\leq i\leq8N$ is a functor, its reduction is an object of $\mathcal P$ of the kind considered, and $\widetilde c_N(4N)=4N$. For $4N>k$ this is a contradiction.\hfill$\square$

So the comparisons between different models of the path object are canonical in $\mathsf{Hot}$, and cannot in general be promoted to isomorphisms in $\mathrm{Cat}\,\mathrm{hot}$, even for a finite ordered set with all arrows marked. The statement was not found in the sources consulted; no priority is claimed.
\end{dcomment}

%%APPEND title="A properness claim which the empty path disproves"
The functor is not a Grothendieck fibration either, already for $X=[1]$: if $c$, of type $\tau$, were a terminal lift of a given arrow of the base, the path that is constant on a first copy of $\tau$ and then follows $c$ would have two different morphisms to $c$, over the two collapses of the doubled type onto $\tau$. The variant in which paths are eventually constant on both sides is smooth and proper, and for $X=[1]$ is neither a fibration nor an opfibration: a path may delay its one step arbitrarily, and no lift is universal. Composing with this variant gives, for every functor $A\to X$, a natural factorization into a weak equivalence followed by a smooth and proper functor whose fibres are the homotopy fibres.

%%FILE III-V
%%APPEND title="A small ind-object that is homotopy invariant sees only 1-types"
The proposition applies to the countable systems of path categories of these pages: for a small relative category, invariance under weak equivalences of all small shapes of the functor $X\mapsto\varinjlim_n\operatorname{Hom}_{\mathrm{Cat}\,\mathrm{hot}}(X,K_n)$ forces it to factor through the fundamental groupoid, and if it is represented by a space, the space is a 1-type. The conclusion depends on allowing shapes $X$ of arbitrary cardinality; for shapes bounded in size relative to the $K_n$ it does not follow.

%%FILE XIII-a
%%APPEND title="The factorization hypothesis seen on three objects"
Corollary~1 of ms.~10 has a consequence for diagrams. Suppose only the factorization of every arrow into a member of $\mathrm{Cof}_W$ followed by a member of $W$, together with an initial object, pushouts and two-out-of-three. Then $\mathrm{Cof}_W\cap W$ is the class of the arrows that remain in $W$ after every pushout (one factorization suffices for this half of the corollary), and $\mathrm{Cof}_W$, $W$ satisfy the axioms of a category with cofibrations in which not every object need be cofibrant; diagrams indexed by finite direct categories can then be replaced by cofibrant ones, so that homotopy pushouts and the other finite homotopy colimits of that kind are defined. The proof is not reproduced here. Two examples limit the statement: in the lattice $\{0<a,b<1\}$ with a suitable $W$ there are enough intrinsic cofibrations and objects that are not cofibrant; and finite sets with the isomorphisms satisfy the finite axioms and not the axioms on infinite compositions and sums.

%%FILE XIIc-a
%%APPEND title="Smoothness and properness depend on a coefficient theory"
For the inclusion $\mathrm{Liss}\cap\mathrm{Prop}\subseteq\mathrm{Fib}_W$ proposed with a ``perhaps'' on ms.~225--226, take for $\mathrm{Liss}$ and $\mathrm{Prop}$ the $W$-smooth and the $W$-proper functors. A $W_\infty$-smooth functor is $W$-smooth and a $W_\infty$-proper one is $W$-proper, for every fundamental localizer $W$, since $W_\infty\subset W$. For $W=W_\infty$ the inclusion holds. For other localizers it fails. The functor from the contractible groupoid with two objects to the group $C_2$ is a fibration and an opfibration, hence smooth and proper, and its pullback along $e\to BC_2$, which is in $W_0$, is the inclusion of two points in a contractible category: so it fails for $W_0$. For the localizer $W_n$ of the $n$-equivalences, $n\geq1$, take a finite category $X$ of the homotopy type of $S^{n+1}$ and the functor $q$ from the category of paths in $X$ issued from a fixed object to $X$ (end point); $e\to X$ is in $W_n$, and its pullback along $q$ is the inclusion of a fibre of the type of $\Omega S^{n+1}$ in a contractible category, which is not in $W_n$. The same path functor over the group $C_2$ gives the failure for rational homology, and a covering with three sheets gives it for the equivalences in $E$-homology, $E$ a nonzero spectrum; these two verifications are not reproduced here.

%%FILE XVI
%%APPEND title="No homology localizer has enough intrinsic fibrations"
The contrast with cofibrations is complete: every fundamental localizer satisfying the relative axiom has enough intrinsic cofibrations in the sense dual to Loc~(7F) as soon as it is accessible, whereas enough intrinsic fibrations force it, with or without accessibility, to be generated by its aspherical objects (the second paragraph of the proof of Proposition~2 at ms.~27).

%%FILE XIIc-b
%%APPEND title="An interval with room at both ends"
For the interval of ms.~228 the inclusion of the two end points, multiplied by a category, need not be an intrinsic $W_\infty$-cofibration, already in the smallest case: there is a sieve inclusion obtained in this way, and a functor from its source to a category $A$, such that the pushout of a contractible category is a category with six objects and twenty arrows having the homology of $S^3$. The inclusion is a weak equivalence; it is its stability under pushout that fails (commentary on ms.~114).

%%FILE II-a
%%APPEND title="Total asphericity is not filteredness"
\textbf{Colimits of categories.}
The property $T(A)$ controls homotopy colimits of diagrams of sets, not strict colimits of diagrams of categories. For a filtered category $A$ and every functor $D:A\to\mathrm{Cat}$, the functor from the total category of $D$ to $\varinjlim D$ is a weak equivalence. For the category $\Delta'$ of the theorem above this fails: let $D([n])$ be the free commutative monoid on $n+1$ generators, regarded as a category with one object. The colimit in $\mathrm{Cat}$ is the monoid $\mathbf N$, of the homotopy type of the circle, whereas the total category has $H_k=\mathbf Z$ for every $k\geq0$: the homology of the stages in degree $k$ is the functor $[n]\mapsto\Lambda^k\mathbf Z^{n+1}$, which is the free abelian group on the representable functor at $[k-1]$, hence projective with colimit $\mathbf Z$. So among nonempty connected categories the filtered ones have the property that all colimits in $\mathrm{Cat}$ are homotopy colimits, the categories with this property satisfy $T$, and the second inclusion is strict. Which categories have the property is not known.
